% GENERATED FILE. Edit canonical lessons, then run npm run generate:books.
% canonical-source-hash: fnv1a64:6cc6a381f603bb35
% canonical-lessons: PA-C309-khabar, PA-C309-hadsa, PA-C309-afvah, PA-C309-raula, PA-C309-marna

\chapter{News from Home}
\label{ch:pa-a2-news-from-home}
\hlchaptermodality{309}

\begin{chapteropening}
\textbf{By the end of this chapter:} \emph{I can pass on a piece of news about something that happened.}

Complete the last lesson of chapter 309: I can pass on a piece of news about something that happened.
\end{chapteropening}

\section[\texorpdfstring{\pa{ਖ਼ਬਰ} (khabar)}{ਖ਼ਬਰ (khabar)}]{\pa{ਖ਼ਬਰ} (khabar) — a piece of news}

\label{lesson:PA-C309-khabar}



\begin{quote}
Before the new one: say the Punjabi for to take out, then the Punjabi for perhaps.
\end{quote}

\subsection*{You'll want to know: \pa{ਖ਼ਬਰ}}
\textbf{\pa{ਖ਼ਬਰ}} — \emph{khabar} — \textquotedblleft{}a piece of news\textquotedblright{}.

\begin{quote}
\pa{ਅਮਨ} \pa{ਨੇ} \pa{ਮੈਨੂੰ} \pa{ਇੱਕ} \pa{ਚੰਗੀ} \pa{ਖ਼ਬਰ} \pa{ਦਿੱਤੀ।}
\end{quote}

\begin{quote}
\emph{Aman ne mainū̃ ikk chaṅgī khabar dittī.}
\end{quote}

\begin{quote}
Aman gave me a piece of good news.
\end{quote}

A feminine noun.

\begin{cousinweb}
\textbf{\pa{ਖ਼ਬਰ}} (\emph{khabar}) is a loan, usually traced through Persian to Arabic \emph{khabar}, \textquotedblleft{}news, a report\textquotedblright{}.
\end{cousinweb}

\begin{grammarlens}[title={What you've built}]
One more word for saying what happened.
\end{grammarlens}

\subsection*{Your turn}
\begin{itemize}
\raggedright
  \item \emph{Say it:} \emph{khabar}
  \item \emph{Say it:} \emph{khabar}, once more
  \item \emph{Say it:} \emph{Aman ne mainū̃ ikk chaṅgī khabar dittī.} — Aman gave me a piece of good news.
  \item \emph{Recall:} say the Punjabi for to take out, then the Punjabi for perhaps, then say \emph{khabar} again
\end{itemize}

\begin{tcolorbox}[breakable,colback=teal!4,colframe=teal!35!black,title={Before you move on}]
What does \pa{ਖ਼ਬਰ} mean? (\textquotedblleft{}A piece of news\textquotedblright{}.) Say it once more.
\end{tcolorbox}
\section[\texorpdfstring{\pa{ਹਾਦਸਾ} (hādsā)}{ਹਾਦਸਾ (hādsā)}]{\pa{ਹਾਦਸਾ} (hādsā) — an accident}

\label{lesson:PA-C309-hadsa}



\begin{quote}
Before the new one: say the Punjabi for perhaps, then the Punjabi for a piece of news.
\end{quote}

\subsection*{You'll want to know: \pa{ਹਾਦਸਾ}}
\textbf{\pa{ਹਾਦਸਾ}} — \emph{hādsā} — \textquotedblleft{}an accident\textquotedblright{}.

\begin{quote}
\pa{ਕੱਲ੍ਹ} \pa{ਸ਼ਾਮ} \pa{ਨੂੰ} \pa{ਸੜਕ} \pa{ਉੱਤੇ} \pa{ਇੱਕ} \pa{ਹਾਦਸਾ} \pa{ਹੋਇਆ।}
\end{quote}

\begin{quote}
\emph{kall shām nū̃ saṛak utte ikk hādsā hoiā.}
\end{quote}

\begin{quote}
There was an accident on the road yesterday evening.
\end{quote}

A masculine noun.

\begin{cousinweb}
\textbf{\pa{ਹਾਦਸਾ}} (\emph{hādsā}) is a loan, usually traced through Persian to Arabic \emph{ḥāditha}, \textquotedblleft{}an event\textquotedblright{}.
\end{cousinweb}

\begin{grammarlens}[title={What you've built}]
One more word for saying what happened.
\end{grammarlens}

\subsection*{Your turn}
\begin{itemize}
\raggedright
  \item \emph{Say it:} \emph{hādsā}
  \item \emph{Say it:} \emph{hādsā}, once more
  \item \emph{Say it:} \emph{kall shām nū̃ saṛak utte ikk hādsā hoiā.} — There was an accident on the road yesterday evening.
  \item \emph{Recall:} say the Punjabi for perhaps, then the Punjabi for a piece of news, then say \emph{hādsā} again
  \item \emph{Write it:} copy \textbf{\pa{ਹਾਦਸਾ}} once, checking the whole word
\end{itemize}

\begin{tcolorbox}[breakable,colback=teal!4,colframe=teal!35!black,title={Before you move on}]
What does \pa{ਹਾਦਸਾ} mean? (\textquotedblleft{}An accident\textquotedblright{}.) Say it once more.
\end{tcolorbox}
\section[\texorpdfstring{\pa{ਅਫ਼ਵਾਹ} (afvāh)}{ਅਫ਼ਵਾਹ (afvāh)}]{\pa{ਅਫ਼ਵਾਹ} (afvāh) — a rumour}

\label{lesson:PA-C309-afvah}



\begin{quote}
Before the new one: say the Punjabi for a piece of news, then the Punjabi for an accident.
\end{quote}

\subsection*{You'll want to know: \pa{ਅਫ਼ਵਾਹ}}
\textbf{\pa{ਅਫ਼ਵਾਹ}} — \emph{afvāh} — \textquotedblleft{}a rumour\textquotedblright{}.

\begin{quote}
\pa{ਬਾਜ਼ਾਰ} \pa{ਵਿੱਚ} \pa{ਹਾਦਸੇ} \pa{ਬਾਰੇ} \pa{ਇੱਕ} \pa{ਅਫ਼ਵਾਹ} \pa{ਸੀ।}
\end{quote}

\begin{quote}
\emph{bāzār vich hādse bāre ikk afvāh sī.}
\end{quote}

\begin{quote}
There was a rumour in the market about an accident.
\end{quote}

A feminine noun.

\begin{cousinweb}
\textbf{\pa{ਅਫ਼ਵਾਹ}} (\emph{afvāh}) is a loan, usually traced through Persian to Arabic \emph{afwāh}, \textquotedblleft{}mouths\textquotedblright{}, the plural of the word for a mouth.
\end{cousinweb}

\begin{grammarlens}[title={What you've built}]
One more word for saying what happened.
\end{grammarlens}

\subsection*{Your turn}
\begin{itemize}
\raggedright
  \item \emph{Say it:} \emph{afvāh}
  \item \emph{Say it:} \emph{afvāh}, once more
  \item \emph{Say it:} \emph{bāzār vich hādse bāre ikk afvāh sī.} — There was a rumour in the market about an accident.
  \item \emph{Recall:} say the Punjabi for a piece of news, then the Punjabi for an accident, then say \emph{afvāh} again
  \item \emph{Write it:} copy \textbf{\pa{ਅਫ਼ਵਾਹ}} once, checking the whole word
\end{itemize}

\begin{tcolorbox}[breakable,colback=teal!4,colframe=teal!35!black,title={Before you move on}]
What does \pa{ਅਫ਼ਵਾਹ} mean? (\textquotedblleft{}A rumour\textquotedblright{}.) Say it once more.
\end{tcolorbox}
\section[\texorpdfstring{\pa{ਰੌਲਾ} (raulā)}{ਰੌਲਾ (raulā)}]{\pa{ਰੌਲਾ} (raulā) — a din; a lot of noise}

\label{lesson:PA-C309-raula}



\begin{quote}
Before the new one: say the Punjabi for an accident, then the Punjabi for a rumour.
\end{quote}

\subsection*{You'll want to know: \pa{ਰੌਲਾ}}
\textbf{\pa{ਰੌਲਾ}} — \emph{raulā} — \textquotedblleft{}a din; a lot of noise\textquotedblright{}.

\begin{quote}
\pa{ਹਾਦਸੇ} \pa{ਤੋਂ} \pa{ਬਾਅਦ} \pa{ਸੜਕ} \pa{ਉੱਤੇ} \pa{ਬਹੁਤ} \pa{ਰੌਲਾ} \pa{ਸੀ।}
\end{quote}

\begin{quote}
\emph{hādse ton bāad saṛak utte bahut raulā sī.}
\end{quote}

\begin{quote}
After the accident there was a lot of noise on the road.
\end{quote}

A masculine noun.

\begin{cousinweb}
\textbf{\pa{ਰੌਲਾ}} (\emph{raulā}) is inherited from Sanskrit \emph{rava}, \textquotedblleft{}a roar, a cry\textquotedblright{}; CDIAL lists Punjabi \emph{raulā}, \textquotedblleft{}noise\textquotedblright{}, under it (CDIAL 10641).
\end{cousinweb}

\begin{grammarlens}[title={What you've built}]
One more word for saying what happened.
\end{grammarlens}

\subsection*{Your turn}
\begin{itemize}
\raggedright
  \item \emph{Say it:} \emph{raulā}
  \item \emph{Say it:} \emph{raulā}, once more
  \item \emph{Say it:} \emph{hādse ton bāad saṛak utte bahut raulā sī.} — After the accident there was a lot of noise on the road.
  \item \emph{Recall:} say the Punjabi for an accident, then the Punjabi for a rumour, then say \emph{raulā} again
  \item \emph{Write it:} copy \textbf{\pa{ਰੌਲਾ}} once, checking the whole word
\end{itemize}

\begin{tcolorbox}[breakable,colback=teal!4,colframe=teal!35!black,title={Before you move on}]
What does \pa{ਰੌਲਾ} mean? (\textquotedblleft{}A din; a lot of noise\textquotedblright{}.) Say it once more.
\end{tcolorbox}
\section[\texorpdfstring{\pa{ਮਰਨਾ} (marnā)}{ਮਰਨਾ (marnā)}]{\pa{ਮਰਨਾ} (marnā) — to die}

\label{lesson:PA-C309-marna}



\begin{quote}
Before the new one: say the Punjabi for a rumour, then the Punjabi for a din.
\end{quote}

\subsection*{You'll want to know: \pa{ਮਰਨਾ}}
\textbf{\pa{ਮਰਨਾ}} — \emph{marnā} — \textquotedblleft{}to die\textquotedblright{}.

\begin{quote}
\pa{ਪਿਛਲੇ} \pa{ਹਫ਼ਤੇ} \pa{ਸਾਡਾ} \pa{ਕੁੱਤਾ} \pa{ਮਰਿਆ} \pa{ਅਤੇ} \pa{ਬੱਚੇ} \pa{ਬਹੁਤ} \pa{ਰੋਏ।}
\end{quote}

\begin{quote}
\emph{pichhle hafte sāḍā kuttā mariā ate bachche bahut roe.}
\end{quote}

\begin{quote}
Our dog died last week, and the children cried a lot.
\end{quote}

Past (it …): \pa{ਉਹ} \pa{ਬਿੱਲੀ} \pa{ਸਰਦੀ} \pa{ਵਿੱਚ} \pa{ਮਰੀ}, \emph{oh billī sardī vich marī}.

\begin{cousinweb}
\textbf{\pa{ਮਰਨਾ}} (\emph{marnā}) is inherited from Sanskrit \emph{marate}, \textquotedblleft{}will die\textquotedblright{}; CDIAL lists Punjabi \emph{marnā}, \textquotedblleft{}dies\textquotedblright{}, under it (CDIAL 9871).
\end{cousinweb}

\begin{grammarlens}[title={What you've built}]
One more word for saying what happened.
\end{grammarlens}

\subsection*{Your turn}
\begin{itemize}
\raggedright
  \item \emph{Say it:} \emph{marnā}
  \item \emph{Say it:} \emph{marnā}, once more
  \item \emph{Say it:} \emph{pichhle hafte sāḍā kuttā mariā ate bachche bahut roe.} — Our dog died last week, and the children cried a lot.
  \item \emph{Recall:} say the Punjabi for a rumour, then the Punjabi for a din, then say \emph{marnā} again
  \item \emph{Write it:} copy \textbf{\pa{ਮਰਨਾ}} once, checking the whole word
\end{itemize}

\begin{tcolorbox}[breakable,colback=teal!4,colframe=teal!35!black,title={Before you move on}]
What does \pa{ਮਰਨਾ} mean? (\textquotedblleft{}To die\textquotedblright{}.) Say it once more.
\end{tcolorbox}
